\SourceLeaf{118}
\Article{24}
The preceding experiments were made at the observatory. The instruments were placed so as to keep iron as far from them as possible. Even so, there is no doubt that the masses of iron scattered throughout the walls, windows and doors of the building, and even the iron parts of the larger astronomical instruments, in which the Earth's magnetic force itself induces magnetism, have a perceptible effect on the suspended needles. The forces arising from this slightly alter both the direction and the intensity of terrestrial magnetism. Our experiments thus give the intensity as modified at the location of instrument $A$, with this disturbance included in the reading. The modification must remain effectively constant as long as the masses of iron stay in place and the elements of terrestrial magnetism itself, namely its intensity and direction, do not change much. Its size remains unknown, though I would scarcely believe it exceeds one or two hundredths of the total value. It should be fairly easy to determine the size at least approximately by experiment. We would observe the simultaneous oscillations of two needles, one suspended at the usual place in the observatory and the other outdoors at a sufficient distance from the building or any other disturbing masses of iron. The needles would then have to exchange places. So far there has been no opportunity to carry out these experiments. The surest remedy, however, will be a special building for magnetic observations, soon to be erected through royal generosity. Iron will be excluded entirely from its construction.

